\chapter{2002年上海卷（秋文）}
\section{填空题}
\begin{problem}
若 \(z \in \mathbf{C}\)，\((3 + z)\i = 1\) (i 为虚数单位)，则 \(z = \fillinblank{}\).
\end{problem}

\begin{problem}
已知向量 \( \bs{a} \) 和 \( \bs{b} \) 的夹角为 \( 120^{\circ} \) ，且 \( |\bs{a}| = 2, |\bs{b}| = 5 \) ，则 \( (2\bs{a} - \bs{b}) \cdot \bs{a} = \)  \fillinblank{}.
\end{problem}

\begin{problem}
方程 \( \log_{8}(1-2\times3^{x})=2x+1 \) 的解 x=\fillinblank{} \fillinblank{}.
\end{problem}

\begin{problem}
若正四棱锥的底面边长为 \( 2\sqrt{3} \) cm，体积为 \( 4 \, cm^{3} \) ，则它的侧面与底面所成的二面的角大小是 \fillinblank{} \fillinblank{}.
\end{problem}

\begin{problem}
在二项式 \((1 + 3x)^2\) 和 \((2x + 5)^3\) 的展开式中，各项系数之和分别记为 \(a_n\)、\(b_n\)，n 是正整数，则 \(\lim_{n \to \infty} \frac{a_n - 2b_n}{3a_n - 4b_n} = \)  \fillinblank{}.
\end{problem}

\begin{problem}
已知圆 \(x^{2} + (y - 1)^{2} = 1\) 和圆外一点 \(P(-2,0)\)，过点 \(P\) 作圆的切线，则两条切线夹角的正切值是 \fillinblank{} \fillinblank{}.
\end{problem}

\begin{problem}
在某次花样滑冰比赛中，主裁判发脚事件，竞赛委员会决定将裁判由原来的9名增至14名，但只任取其中7名裁判的评分作为有效分，若14名裁判中有2人受脚，则有效分中没有受脚裁判的评分的概率是 \fillinblank{}. （结果用数值表示）
\end{problem}

\begin{problem}
抛物线 \((y - 1)^2 = 4(x - 1)\) 的焦点坐标是 \fillinblank{} \fillinblank{}.
\end{problem}

\begin{problem}
某工程由下列工序组成，则工程总时数为 天. 工序 a b c d e f 紧前工序 - -a, b c c d, e 工时数(天) 2 3 2 5 4 1
\end{problem}

\begin{problem}
设函数 \( f(x) = \sin 2x \)，若 \( f(x + t) \) 是偶函数，则 \( t \) 的一个可能值是 \fillinblank{} \fillinblank{}.
\end{problem}

\begin{problem}
若数列 \(\{a_{n}\}\) 中，\(a_{1} = 3\)，且 \(a_{n+1} = a_{n}^{2}\)（\(n\) 是正整数），则数列的通项 \(a_{n} = \)  \fillinblank{}.
\end{problem}

\begin{problem}
已知函数 \( y = f(x) \) （定义域为 \( D \)，值域为 \( A \)）有反函数 \( y = f^{-1}(x) \)，则方程 \( f(x) = 0 \) 有解 \( x = a \)，且 \( f(x) > x (x \in D) \) 的充要条件是 \( y = f^{-1}(x) \) 满足  \fillinblank{}.
\end{problem}

\section{单选题}
\begin{problem}
如图，与复平面中的阴影部分 (含边界) 对应的复数量集合是（\quad）

\choices
    {\(\left\{|z| = 1, \operatorname{Re} z \geqslant \frac{1}{2}: z \in \mathbf{C}\right\}\)}
    {\(\left\{|z| \leqslant 1, \operatorname{Re} z \geqslant \frac{1}{2}: z \in \mathbf{C}\right\}\)}
    {\(\left\{z\mid |z| = 1,\operatorname {Im}z\geqslant \frac{1}{2},z\in \mathbf{C}\right\}\)}
    {\(\left\{z\mid |z|\leqslant 1,\operatorname {Im}z\geqslant \frac{1}{2},z\in \mathbf{C}\right\}\)}
\end{problem}

\begin{problem}
已知直线 \(l\)、\(m\)、平面 \(\alpha\)、\(\beta\)，且 \(l \perp \alpha, m \subset \beta\)，给出下列四个命题. \circled{1} 若 \(\alpha // \beta, l \perp m\); \circled{2} \(l \perp m, \alpha // \beta\); \circled{3} 若 \(\alpha \perp \beta\)，则 \(l // m\); \circled{4} 若 \(l // m\)，\(\alpha \perp \beta\). 其中正确命题的个数是（\quad）

\choices
    {1 个}
    {2 个}
    {3 个}
    {4 个}
\end{problem}

\begin{problem}
函数 \( y = x + \sin |x| \)，\( x \in [-\pi, \pi] \) 的大致图象是（\quad）

\bitmapfigure[max width=0.9\linewidth,max totalheight=4.0cm]{img/2002/shanghai_liberal/q15_options.png}
\end{problem}

\begin{problem}
一般地，家庭用电量 (千瓦时) 与气温 ( \( ^{\circ} \) C) 有一定的关系. 图 (1) 表示某年 12 个月中每月的平均气温，图 (2) 表示某家庭在这 12 个月中每月的用电量，根据这些信息，以下关于该家庭用电量与气温间关系的叙述中，正确的是（\quad）

\bitmapfigure[max width=0.86\linewidth,max totalheight=4.8cm]{img/2002/shanghai_liberal/q16_fig1.png}

\choices
    {气温最高时，用电量最多}
    {气温最低时，用电量最少}
    {当气温大于某一值时，用电量随气温增高而增加}
    {当气温小于某一值时，用电量随气温降低而增加}
\end{problem}

\section{解答题}
\begin{problem}
如图，在直三棱柱 \(ABO - A'B'O'\) 中，\(OO' = 4\)，\(OA = 4\)，\(OB = 3\)，\(\angle AOB = 90^\circ\)，\(D\) 是线段 \(A'B'\) 的中点，\(P\) 是侧棱 \(BB'\) 上的一点，若 \(OP \perp BD\)，求 \(OP\) 与底面 \(AOB\) 所成角的大小. (结果用反三角函数值表示)
\end{problem}

\begin{problem}
已知点 \(A(-\sqrt{3},0)\) 和 \(B(\sqrt{3},0)\)，动点 \(C\) 到 \(A\)、\(B\) 两点的距离之差的绝对值为 2，点 \(C\) 的轨迹与直线 \(y = x - 2\) 交于 \(D\)、\(E\) 两点，求线段 \(DE\) 的长.
\end{problem}

\begin{problem}
已知函数 \( f(x) = x^2 + 2ax + 2, x \in [-5, 5] \). (1) 当 \(a = -1\) 时，求函数 \(f(x)\) 的最大值与最小值. (2) 求实数 a 的取值范围，使 \( y = f(x) \) 在区间 \( [-5, 5] \) 上是单调函数.
\end{problem}

\begin{problem}
某商场在促销期间规定：商场内所有商品按标价的 \(80\%\) 出售，同时，当顾客在该商场内消费满一定金额后，按如下方案获得相应金额的奖券： 消费金额的范围 200,400) 400,500) 500,700) 700,900) ... 获得奖券的金额 30 60 100 130 ... 根据上述促销方法，顾客在该商场购物可以获得双重优惠，例如：购买标价为400元的商品，则消费金额为320元. 获得的优惠额为： \( 400 \times 0.2 + 30 = \) 购买商品获得的优惠额 110(元)，设购买商品得到的优惠率 \(= \frac{\text{商品的标价}}{\text{商品的标价}}\)，试问： (1) 若购买一件标价为 1000 元的商品，顾客得到的优惠率是多少\(\perp\) (2) 对于标价在 [300, 800] (元) 内的商品，顾客购买标价为多少元的商品，可得到不小于 \( \frac{1}{4} \) 的优惠率\(\perp\)
\end{problem}

\begin{problem}
已知函数 \( f(x) = a \times b^x \) 的图象过点 \( A\left(4, \frac{1}{4}\right) \) 和 \( B(5, 1) \). (1) 求函数 \( f(x) \) 的解析式; (2) 记 \(a_{n} = b_{n} g_{n} f(n), n\) 是正整数，\(S_{n}\) 是数列 \(\{a_{n}\}\) 的前 \(n\) 项和，解决关于 \(n\) 的不等式 \(a_{n} S_{n} \leqslant 0\); (3) 对于 (2) 中的 \(a_{n}\) 与 \(S_{n}\) 互数，整 \(90\%\) 是否数列 \(\{a_{n} S_{n}\}\) 中的项\(\perp\) 若是，则求出相应的项数; 若不是，则说明理由.
\end{problem}

\begin{problem}
规定 \( C_{n}^{m} = \frac{x(x - 1)\cdots(x - m + 1)}{n!} \)，其中 \( x \in \R, m \) 是正整数，且 \( C_{n}^{0} = 1 \)，这是组合数 \( C_{n}^{m}(n, m) \) 是正整数，且 \( m \leqslant n \) 的一种推广. (1) 求 \( C_{n+1}^{3} \) 的值. (2) 设 \(x > 0\)，当 \(x\) 为何值时，\(\frac{C_2^1}{(C_2^1)^2}\) 取得最小值\(\perp\) (3) 组合数的两个性质: \circled{1} \(C_{n}^{m} = C_{n}^{n - m}\); \circled{2} \(C_{n}^{m} + C_{n}^{m - 1} = C_{n + 1}^{m}\)，是否都能推广到 \(C_{n}^{m}\) (\(x \in \R, m\) 是正整数) 的情形\(\perp\) 若能推广，则写出推广的形式并给出证明; 若不能，则说明理由.
\end{problem}

